
\documentclass[11pt,a4paper]{article}

\usepackage{amsmath,amssymb,amsthm}
\usepackage{geometry}
\geometry{margin=1in}
\usepackage{hyperref}
\usepackage{enumitem}

\newtheorem{theorem}{Theorem}[section]
\newtheorem{lemma}[theorem]{Lemma}
\newtheorem{proposition}[theorem]{Proposition}
\newtheorem{corollary}[theorem]{Corollary}

\theoremstyle{definition}
\newtheorem{definition}[theorem]{Definition}
\newtheorem{remark}[theorem]{Remark}
\newtheorem{observation}[theorem]{Observation}

\newcommand{\MQ}{\mathrm{MQ}}
\newcommand{\FF}{\mathbb{F}}
\newcommand{\Fix}{\mathrm{Fix}}
\newcommand{\RM}{\mathrm{RM}}
\newcommand{\cl}{\mathrm{cl}}
\newcommand{\ev}{\mathrm{ev}}
\DeclareMathOperator{\supp}{supp}
\DeclareMathOperator{\wt}{wt}
\DeclareMathOperator{\artanh}{artanh}
\DeclareMathOperator{\Rank}{Rank}

\title{\textbf{A Note on Binary Quadratic Systems and their relation to complexity theory}}
\author{%
  Gabriele Radici\thanks{Department of Mathematics, University of Trento, Italy.
  \texttt{gabriele.radici@studenti.unitn.it}}
  \and
  Massimiliano Sala\thanks{Department of Mathematics, University of Trento, Italy.
  \texttt{maxsalacodes@gmail.com}}%
}
%\date{August 10, 2026}

\begin{document}

\maketitle


\begin{abstract}
Deciding whether a system of multivariate quadratic equations over
$\mathbb F_2$ has a solution is a classical NP-complete problem, and remains so for
square systems, with as many equations as variables. 
The hardness of this problem is one of the
cornerstones of nowadays post-quantum cryptography.
Let $\MQ_0(n)$ and $\MQ_1(n)$ denote the sets of square quadratic
systems in $n$ variables having respectively no solutions and exactly
one solution. $\cup_{n\geq 2} \MQ_0(n)$ is a coNP-complete language, while
$\cup_{n\geq 2} \MQ_1(n)$ lies in DP. 
It is known that
$\lim_{n\to \infty} |\MQ_1(n)|/|\MQ_0(n)|=1$.
Here we prove the explicit finite-$n$ bounds
\[
|\MQ_0(n)|<|\MQ_1(n)|
\le
\left(1+\frac{1}{2^n-1}\right)|\MQ_0(n)|,
\]

We find of interest that the two previous inequalities are equivalent to the existence of some injections
$
\MQ_0(n)\hookrightarrow\MQ_1(n)
$
(whose images omit at most a $2^{-n}$ fraction of $\MQ_1(n)$).
Whether such an injection can be given explicitly, or computed or inverted in
polynomial time, is open.


More generally, let $Q_d$ be the space of polynomial functions
$(\FF_2)^n\to\mathbb F_2$ of degree at most $d$, and let $\alpha_k$
count square systems in $(Q_d)^n$ having exactly $k$ solutions. Then
\[
\alpha_0<\alpha_1
\le
\left(1+\frac{1}{2^n-1}\right)\alpha_0\,, \qquad 2\le d\le n \,.
\]
The proof combines matroid and coding-theoretic methods. We interpret
$(\FF_2)^n$ as the ground set of the evaluation matroid of $Q_d$,
express $\alpha_0$ and $\alpha_1$ through characteristic polynomials,
and use a Whitney-type sign-reversing involution to show that the only
terms that can push $\alpha_1-\alpha_0$ below $\alpha_1/2^n$ come from the
elements of a matroid port. These are identified with
minimal-support words of the Reed--Muller code
$\RM(n-d-1,n)=\RM(d,n)^\perp$; the required estimate then follows
from the MacWilliams identity, the minimum-distance bound $2^{d+1}$,
and the even-weight structure of the code.
\end{abstract}


\section{Introduction}\label{sec:intro}

In this paper we will present results of an algebraic, geometric and combinatorial nature,
but which may open the path to results in complexity theory. These results are also connected to post-quantum cryptography.
The reader unfamiliar with complexity theory will find in Section~\ref{sec:complexity} a collection of notions and statements useful to understand the connection with our algebraic results.

Let $\FF$ denote the finite field with $2$ elements, usually denoted
by $\FF_2$ or $\mathrm{GF}(2)$.
Deciding whether a system of multivariate quadratic equations over
$\FF$ has (at least) a solution is a classical NP-complete problem
\cite{fraenkelyesha,gareyjohnson}. As shown in Section~\ref{sec:complexity}, NP-completeness persists when one restricts to
\emph{square} systems, that is, to systems with as many equations as variables.
Beyond its complexity-theoretic status, the hardness of this problem (in the more general finite-field setting) is the main hardness assumption underlying
multivariate post-quantum cryptography. For example, of the nine candidates advanced by NIST to
the third round of its additional digital signature process in May 2026, four are variants of the Unbalanced Oil and Vinegar scheme \cite{uov}: UOV itself, MAYO \cite{mayo}, QR-UOV \cite{qruov} and SNOVA \cite{snova}, whose security rests on the difficulty of solving systems of quadratic equations over a finite field
\cite{nistir8610}. 

For $n\ge2$, let $\MQ_0(n)$ and $\MQ_1(n)$ denote the sets of square
quadratic systems in $n$ variables over $\FF$ having respectively
no solutions and exactly one solution. The language
$
\bigcup_{n\ge2}\MQ_0(n)
$
is coNP-complete, whereas
$\bigcup_{n\ge2}\MQ_1(n)$
belongs to DP, not known to be either in NP or coNP. 
Moreover, $\bigcup_{n\ge2}\MQ_1(n)$ is coNP-hard and, under randomized
reductions, NP-hard. For the unique satisfiability problem these are the theorems of Blass and Gurevich \cite{blassgurevich} and of Valiant
and Vazirani \cite{valiantvazirani}; their transfer to square quadratic systems is discussed in Section~\ref{sec:complexity}.

The
two families have remarkably close cardinalities. It is known
\cite{fuscobach} that
\[
\lim_{n\to\infty}
\frac{|\MQ_1(n)|}{|\MQ_0(n)|}=1.
\]
Here we prove the explicit finite-$n$ bounds
\begin{equation}\label{eq:intro-mq}
|\MQ_0(n)|<|\MQ_1(n)|
\le
\left(1+\frac{1}{2^n-1}\right)|\MQ_0(n)|.
\end{equation}
Therefore, with only a relative excess of order $2^{-n}$, a
uniformly-random square quadratic system is 
\underline{more likely to have exactly
one solution than none}.\\


The two inequalities in
\eqref{eq:intro-mq} are equivalent to the existence of an
injection
\[
\MQ_0(n)\hookrightarrow\MQ_1(n)\,,\qquad \mbox{for any } n\geq 2\,,
\]
whose image omits at most a $2^{-n}$ fraction of $\MQ_1(n)$.
Whether such an injection can be given explicitly, or computed/inverted
in polynomial time, is open: an injection of that kind would relate a
coNP-complete language to a language in DP. While we make no
complexity-theoretic claim in this paper,
in Section~\ref{sec:complexity} we will discuss some consequences.
\\

Our result for quadratic systems is a special case of the more general
statement that we prove in this paper. Let
$n\ge 2$, $V=\FF^n$, $N=|V|=2^n$ and write
$$
Q:=\FF[x_1,\dots,x_n]/\langle x_i^2-x_i\rangle
$$
for the ring of polynomial functions $V\to\FF$. 
For $0\le d\le n$,
let $Q_d\subseteq Q$ be the subspace of functions of degree at most
$d$, of dimension
$
D=D(n,d)=\sum_{j=0}^{d}\binom n j
$.\\
A \emph{square system of degree $\le d$} is an $n$-tuple
$F\in (Q_d)^{\,n}$ with variety $X(F)\subseteq V$, that is,
\[
F=(f_1,\dots,f_n)\in (Q_d)^{\,n}, \qquad 
X(F):=\{v\in V \mid f_1(v)=\dots=f_n(v)=0\}.
\]
For $0\le k\le N$, set
\[
\alpha_k=\alpha_k(n,d)
:=
\bigl|\{F\in (Q_d)^{\,n}:|X(F)|=k\}\bigr|.
\]
Thus, for $d=2$,
\[
|\MQ_k(n)|=\alpha_k(n,2).
\]

\newpage

Our main result is the following.
%
\begin{theorem}[restated with proof as Theorem~\ref{thm:main}]\label{thm:intro}
Let $n\ge 2$ and $2\le d\le n$. Then
\begin{enumerate}[label=\textnormal{(\roman*)}]
\item
\[
\alpha_1-\alpha_0\le\frac{\alpha_1}{N},
\]
with equality if and only if $d=n$;

\item
\[
\alpha_1>\alpha_0;
\]

\item if $d\le n-2$, then
\[
\frac{\alpha_1}{N}-\frac{N^{D-3}}{3}
\le
\alpha_1-\alpha_0
\le
\frac{\alpha_1}{N},
\qquad
\frac{\alpha_1}{N}
\ge
N^{D-4}\frac{(N-1)(N-2)(2N+3)}{6}
\ge
\frac{N^{D-1}}{4};
\]
hence
\[
\left(1-\frac4{3N^2}\right)\frac{\alpha_1}{N}
\le
\alpha_1-\alpha_0
\le
\frac{\alpha_1}{N}.
\]
\end{enumerate}
\end{theorem}

Theorem~\ref{thm:intro} specializes for quadratic systems to
\eqref{eq:intro-mq}, since
\[
\alpha_1-\alpha_0\le\frac{\alpha_1}{N}
\quad\Longleftrightarrow\quad
\alpha_1\le
\left(1+\frac1{N-1}\right)\alpha_0\,.
\]
The equality in its upper bound occurs if and only if
$n=2$.

Two boundary cases of Theorem~\ref{thm:intro} are classical. For
$d=n$, $Q_n$ is the space of all functions $V\to\FF$, the number of
solutions of a uniformly random square system is exactly binomial with
parameters $N$ and $1/N$ \cite[Theorem~1]{jain}, and (i) holds with
equality; by Theorem~\ref{thm:intro}(i) itself, $d=n$ is the only
degree $\ge2$ for which the number of solutions is binomial (compare
\cite[Corollary~3.6]{jain}: the events ``$F$ vanishes at $v$'', $v\in
V$, are independent only in that case), which is what makes the case
$d<n$ non-trivial. For $d=1$ the number of solutions of a random
square linear system over $\FF$ has the classical distribution
governed by the rank of a random matrix \cite[Chapter~3]{kolchin},
and there the inequality goes the other way: $\alpha_0>\alpha_1$ for
every $n\ge3$ (Remark~\ref{rem:final}). So the hypothesis $d\ge2$ in
Theorem~\ref{thm:intro} cannot be dropped.
\\

We briefly describe the proof. Section~\ref{sec:prelim} collects the facts on matroids that we use. In Section~\ref{sec:matroid}, the point set $V$ is made
into a matroid $M$ by declaring $S\subseteq V$
independent when the evaluation vectors of its points on $Q_d$ are
linearly independent. Inclusion--exclusion then gives
\[
\alpha_0=\chi_M(N),
\qquad
\alpha_1=N\chi_{M/v}(N),
\]
where $\chi$ denotes the characteristic polynomial and $v\in V$ is
arbitrary.

In Section~\ref{sec:reduction}, a sign-reversing involution in the
spirit of Whitney's broken-circuit theorem, adapted to the port of $M$
at $v$, shows that the only terms that can push
$\alpha_1-\alpha_0$ below $\alpha_1/N$ come from the minimal sets of
the port of odd cardinality. More precisely, if $\mu_j$ denotes the
number of minimal sets of the port of cardinality $j$, we obtain
\[
\alpha_1-\alpha_0
\ge
\frac{\alpha_1}{N}
-
\sum_{j\ {\rm odd}}\mu_jN^{D+1-j}.
\]

In Section~\ref{subsec:rm-codes}, these minimal sets are identified
with the minimal-support words through $v$ of the Reed--Muller code
\[
\RM(n-d-1,n)=\RM(d,n)^\perp,
\]
so that the $\mu_j$ become weight-enumerator coefficients. In
Section~\ref{sec:rm-bound} the resulting error term is controlled through the
MacWilliams identity: the minimum distance $2^{d+1}\ge8$ implies agreement with
the corresponding binomial model through the first seven moments, while the
even-weight structure gives the required sign for the upper bound, and a
sixth-moment estimate makes the error term small enough.
Section~\ref{sec:mainthm} assembles these ingredients into a proof of
Theorem~\ref{thm:intro}, and shows that its hypothesis $d\ge2$ cannot be
dropped.

Finally, in Section~\ref{sec:complexity} some comments are given on the
relation between the research carried out in this paper and complexity
theory.


\section{Preliminaries on matroids}\label{sec:prelim}

Matroids, introduced by Whitney \cite{whitney35}, abstract linear dependence; we refer to \cite{welsh,oxley} for the basic theory. We use the rank-function axiomatization.

\begin{definition}\label{def:matroid-rank}
A \emph{matroid} $M$ is a pair $(E,r)$, where $E$ is a finite set, the \emph{ground set}, and $r:2^E\to\mathbb Z_{\ge0}$ satisfies, for all $A,B\subseteq E$:
\begin{enumerate}[label=\textnormal{(R\arabic*)}]
    \item $0\le r(A)\le|A|$;
    \item if $A\subseteq B$ then $r(A)\le r(B)$;
    \item $r(A\cup B)+r(A\cap B)\le r(A)+r(B)$.
\end{enumerate}
The rank of $M$ is $r(M):=r(E)$.
\end{definition}

Matroids are also usually defined by axioms on the independent sets; the two
axiomatizations are equivalent, in the following sense.

\begin{theorem}[\cite{whitney35,welsh}]\label{thm:rank-equivalence}
Call $\mathcal I\subseteq2^E$ an \emph{independence system} if
\begin{itemize}
    \item $\emptyset\in\mathcal I$; \item subsets of members of $\mathcal I$ are in $\mathcal I$; \item if $I,J\in\mathcal I$ and $|I|<|J|$ there is $e\in J\setminus I$ with $I\cup\{e\}\in\mathcal I$.
\end{itemize}
The assignments
\[
r\ \longmapsto\ \mathcal I_r:=\{I\subseteq E: r(I)=|I|\},
\qquad
\mathcal I\ \longmapsto\ r_{\mathcal I}(A):=\max\{|I| : I\subseteq A,\ I\in\mathcal I\}
\]
are mutually inverse bijections between the functions $r:2^E\to\mathbb Z_{\ge0}$
satisfying \textnormal{(R1)}--\textnormal{(R3)} and the independence systems on $E$.
\end{theorem}

\begin{remark}
The correspondence must be stated in this form. It is \emph{not} true that an
arbitrary $r:2^E\to\mathbb Z_{\ge0}$ satisfies \textnormal{(R1)}--\textnormal{(R3)}
as soon as $\mathcal I_r$ is an independence system: for $E=\{1,2\}$ and
$r(\emptyset)=0$, $r(\{1\})=r(\{2\})=1$, $r(\{1,2\})=5$ one gets
$\mathcal I_r=\{\emptyset,\{1\},\{2\}\}$, which satisfies the three axioms, while
$r$ violates \textnormal{(R1)}. What recovers $r$ from $\mathcal I_r$ is the
displayed formula $r=r_{\mathcal I_r}$, valid for rank functions.
\end{remark}
\textbf{Notation:} If not better specified, throughout all the paper $M$ will denote a matroid $(E,r)$. By some abuse of notation, we will write $S \in M$, if $S \in 2^E$, and $T \subseteq M$, if $T \subseteq 2^E$. \\

From \cite{welsh} and \cite{whitney35}, we also bring the following definition and basic facts on matroids.

\begin{definition}\label{def:basic}
Let $M = (E, r)$ be a matroid with rank function $r$. We define:
\begin{enumerate}
    \item \textbf{Independent sets}: $\mathcal I:=\{I\subseteq E: r(I)=|I|\}$; the other subsets are \textbf{dependent}.
    \item \textbf{Bases}: maximal independent sets. 
    \item \textbf{Circuits} $\mathcal C(M)$: minimal dependent sets.
    \item \textbf{Loop}: $e\in E$ with $r(\{e\})=0$. 
    \item \textbf{Coloop}: $e\in E$ with $r(E\setminus\{e\})=r(E)-1$. \item Two non-loops $e\ne f$ are \textbf{parallel} if $r(\{e,f\})=1$.
    \item \textbf{Closure}: $\cl:2^E\to2^E$, $\cl(A):=\{e\in E: r(A\cup\{e\})=r(A)\}$.
    \item \textbf{Deletion}: for $T\subseteq E$, $M\setminus T$ is the matroid on $E\setminus T$ with rank $r_{M\setminus T}(A)=r(A)$.
    \item \textbf{Contraction}: for $T\subseteq E$, $M/T$ is the matroid on $E\setminus T$ with rank $r_{M/T}(A)=r(A\cup T)-r(T)$.
    \item \textbf{Characteristic polynomial}: $\chi_M(\lambda)=\sum_{A\subseteq E}(-1)^{|A|}\lambda^{r(M)-r(A)}$.
\end{enumerate}
\end{definition}

\begin{remark}\label{rem:facts}
The following are standard \cite{welsh,oxley}.
\begin{enumerate}[label=\textnormal{(\alph*)}]
    \item All bases have cardinality $r(M)$.
    \item A subset $C \subseteq E$ is a circuit iff $r(C) = |C| - 1$ and $r(C \setminus \{e\}) = |C| - 1$ for all $e \in C$.
    \item The closure operator $\cl: 2^E \to 2^E$ satisfies the following properties. For any set $S \subseteq E$: 
    \begin{itemize}
        \item Extensivity: $S \subseteq \cl(S)$;
        \item Monotonicity: if $S \subseteq S'$ then $\cl(S) \subseteq \cl(S')$;
        \item Idempotency: $\cl(\cl(S)) = \cl(S)$.
    \end{itemize}
    \item If $I$ is independent and $e\in\cl(I)\setminus I$, then $I\cup\{e\}$ contains exactly one circuit, and that circuit contains $e$. It is called \textbf{fundamental circuit}, we denote it as $I^e$.
    \item $r_{M/T}$ is still a rank function. We have that, since $A \cap T = \emptyset$ and submodularity of $r$, $r(A \cup T) - r(T) \le |A|$. To prove submodularity of $r_{M/T}$, one notices that $r(A \cup B \cup T) + r((A \cap B) \cup T) = r((A \cup T) \cup (B \cup T)) + r((A \cup T) \cap (B \cup T)) \le r(A \cup T) + r(B \cup T)$.
    \item The characteristic polynomial of $M\setminus T$ is $\chi_{M\setminus T} = \sum_{S \subseteq E \setminus T}(-1)^{|S|}\lambda^{r(M \setminus T) - r(S)}$;
    \item The characteristic polynomial of $M/T$ is $\chi_{M/T} = \sum_{T \subseteq S \subseteq E}(-1)^{|S| - |T|}\lambda^{r(M) - r(S)} =$ \\ $ = \sum_{S \subseteq E \setminus T}(-1)^{|S|}\lambda^{r(M) - r(S \cup T)}$
\end{enumerate}
\end{remark}

Characteristic polynomials of deletion and contraction matroids provide a useful identity:

\begin{proposition}[\cite{welsh,oxley}]\label{identity}
If $e\in E$ is neither a loop nor a coloop, then \[ \chi_M(\lambda)=\chi_{M\setminus e}(\lambda)-\chi_{M/e}(\lambda)  \,.
\]
\end{proposition}

\section{Boolean polynomial systems as matroids} \label{sec:matroid}

The point set $V=\FF^n$ carries a matroid structure in which a set of points is
independent when the corresponding evaluation functionals on $Q_d$ are linearly
independent. The purpose of this section is to set that structure up and to record
its one consequence that the rest of the paper uses: the numbers $\alpha_0$ and
$\alpha_1$ of systems with no solution and with exactly one solution are values of
the characteristic polynomial of this matroid and of one of its contractions.

\subsection{Rank function}

Let $\mathbb{F} = \mathbb{F}_2$ be the field with two elements, $V = \mathbb{F}^n$ be the $n$-dimensional vector space over $\mathbb{F}$, and $N = 2^n = |V|$.
Set $Q := \mathbb{F}[x_1, \dots, x_n] / \langle x_i^2 - x_i : i = 1, \dots, n \rangle$ and, for $0 \le d \le n$, let $Q_d$ denote the space of polynomials in $Q$ up to degree $d$. The dimension of $Q_d$ is 
\[
D = \sum_{j=0}^d \binom{n}{j}.
\]
Let $M_d$ be the set of monomials of $Q_d$. Since $M_d$ is a basis for $Q_d$, $Q_d$ is isomorphic to $\mathbb{F}^D$. 
We order monomials in $M_d$ with respect to the \textit{graded lexicographic} ordering. Namely, let $1<x_1<x_2<\dots<x_n$, then $x_1^{i_1}\dots x_n^{i_n} < x_1^{j_1}\dots x_n^{j_n}$ (with $i_k,j_k \in \{0,1\}$) if and only if: \begin{itemize}
    \item $\sum_{k=1}^ni_k < \sum_{k=1}^nj_k$;
    \item $\sum_{k=1}^ni_k = \sum_{k=1}^nj_k$ and $0 = j_{\overline{k}} < i_{\overline{k}} = 1$, where $\overline{k} := \min\{k \ | \ i_k \neq j_k \}$.
\end{itemize}
Define the evaluation map \[ \mathrm{ev} : V \to \mathbb{F}^D,\  \mathrm{ev}(v) := (m(v) : m \in M_d).\] Here $\mathbb{F}^D$ carries two roles, identified throughout by the standard scalar product: it is the space of coefficient vectors of the polynomials of $Q_d$, and, through $\overline{f}\mapsto\overline{f}\cdot\mathrm{ev}(v)$, it is also the space of linear functionals on $Q_d$; the vector $\mathrm{ev}(v)$ is the functional ``evaluate at $v$''. Extending $\mathrm{ev}$ to the power set of $V$ by $\mathrm{ev}(S) = \{\mathrm{ev}(v) : v \in S\}$ for $S \subseteq V$, we define the set function:
\[
\overline{r}(S) := \mathrm{Rank}(\mathrm{ev}(S)), \quad \text{for } S \subseteq V.
\]

From now on, we assume $d \ge 2$. So,
\begin{remark}
    $\mathrm{ev}$ is injective, indeed $\mathrm{ev}(v) = (1, v, (m(v))_{m \in M_d, \deg m \ge 2})$. Hence $\mathrm{ev}(v)=\mathrm{ev}(w)$ implies $v = w$.
\end{remark}

\begin{lemma}\label{obs:boolean-matroid}
The function $S \mapsto \overline{r}(S)$ is a matroid rank function on ground set $V$.
\end{lemma}

\begin{proof}
We directly verify axioms (R1)--(R3) of Definition~\ref{def:matroid-rank} for $\overline{r}(S)$:

\textbf{(R1)}: $\mathrm{ev}(S)$ is a collection of $|S|$ vectors in $\mathbb{F}^D$. The dimension of $\mathrm{span}(\mathrm{ev}(S))$ cannot exceed the number of vectors $|S|$ nor can it be negative, so $0 \le \overline{r}(S) \le |S|$.

\textbf{(R2)}: If $A \subseteq B \subseteq V$, then $\mathrm{ev}(A) \subseteq \mathrm{ev}(B)$. Taking linear spans yields $\mathrm{span}(\mathrm{ev}(A)) \subseteq \mathrm{span}(\mathrm{ev}(B))$, so $\dim \mathrm{span}(\mathrm{ev}(A)) \le \dim \mathrm{span}(\mathrm{ev}(B))$, i.e., $\overline{r}(A) \le \overline{r}(B)$.

\textbf{(R3)}: Let $A, B \subseteq V$. Set $U = \mathrm{span}(\mathrm{ev}(A))$ and $W = \mathrm{span}(\mathrm{ev}(B))$. Then $U + W = \mathrm{span}(\mathrm{ev}(A \cup B))$. Since $\mathrm{ev}(A \cap B) \subseteq \mathrm{ev}(A) \cap \mathrm{ev}(B) \subseteq U \cap W$, we have $\mathrm{span}(\mathrm{ev}(A \cap B)) \subseteq U \cap W$, so $\overline{r}(A \cap B) \le \dim(U \cap W)$. Applying the standard dimension formula for vector subspaces:
\begin{align*}
\overline{r}(A \cup B) + \overline{r}({A \cap B}) &= \dim(U + W) + \overline{r}(A \cap B) \\
&\le \dim(U + W) + \dim(U \cap W) \\
&= \dim U + \dim W = \overline{r}(A) + \overline{r}(B). \qedhere
\end{align*}
\end{proof} 

\begin{corollary}
    $(V,\overline{r})$ is a matroid. We call it \textbf{evaluation matroid} (of degree $d$)
\end{corollary}

\begin{remark}
    $\mathrm{ev}(V)$ spans the whole $\FF^D$ (since no non-zero polynomial vanishes on all $V$) hence, $\overline{r}(V) = \dim\FF^D = D$. 
\end{remark}

\textbf{Notation:} Established that $\overline{r}$ is a rank function on the powerset of $V$, for compactness we will write $r_S := \overline{r}(S)$, $S \subseteq V$. \\

\subsection{Translation of matroid concepts}

Having proven that $(V, \overline{r})$ is a matroid, we specialize general definitions in our particular case for better understanding.

\begin{enumerate}
    \item \textbf{Independent Sets $\mathcal{I}$}: Point sets $I \subseteq V$ whose monomial evaluation vectors $\mathrm{ev}(I)$ are linearly independent in $\mathbb{F}^D$.
    \item \textbf{Bases $\mathcal{B}(M)$}: Maximal point sets in $V$ whose evaluation vectors form a vector space basis for $\mathrm{span}(\mathrm{ev}(V)) = \mathbb{F}^D$.
    \item \textbf{Circuits $\mathcal{C}(M)$}: Minimal point sets $C \subseteq V$ exhibiting a non-trivial linear dependency $\sum_{v \in C} c_v \mathrm{ev}(v) = 0$ over $\mathbb{F}$. As shown in Section \ref{subsec:rm-codes}, these correspond to minimal support non-zero codewords in the dual Reed-Muller code $\mathrm{RM}(d, n)^\perp = \mathrm{RM}(n - (d + 1), n)$.
    \item \textbf{Deletion $M \setminus T$}: Restricting evaluation to the subset of points $V \setminus T$.
    \item \textbf{Contraction $M/T$}: Projecting the coefficient space $\mathbb{F}^D$ modulo the subspace spanned by $\mathrm{ev}(T)$.
    \item \textbf{Characteristic Polynomial $\chi_M(\lambda)$}: $\sum_{S \subseteq V} (-1)^{|S|} \lambda^{D - r_S}$.
\end{enumerate}

Concerning loops, coloops and parallel elements in the evaluation matroid, we prove that 

\begin{lemma}\label{lem:coloop}
$M$ has no loops, and no two points are parallel. If $2 \le d\le n-1$ then $M$ has no coloops; if $d=n$ every point is a coloop.
\end{lemma}

\begin{proof}
For any $v \in V$, $\ev(v)$ has first coordinate $1(v)=1$, so it is non-zero, hence it cannot have zero rank. \\ $v,w \in V$ are parallel iff $\ev(v)=\ev(w)$, impossible since $\mathrm{ev}$ is injective. \\
If $2 \le d\le n-1$ then for any $v \in V$, $\{f\in Q_d: f(w)=0\ \forall w\in V \setminus \{v\}\}= \{0\}$, because the only function vanishing on $V\setminus\{v\}$ and not on $v$ is the indicator of $v$, of degree $n$; hence $r(V\setminus\{v\})=D=r_V$. If $d=n$ then $Q_d=Q$ and $D=N$, so $\ev(V)$ is a basis of $\FF^N$ and every point is a coloop.
\end{proof}



\subsection{Counting systems with 0 and 1 solutions}

For $S\subseteq V$ let
\[
H_S:=\{f\in Q_d: f(v)=0\ \forall v\in S\}
\]
Let $\overline{f} \in \mathbb{F}^D$ represent the vector of coefficients of the polynomial $f = \sum f_i m_i$, by construction $f(v) = \overline{f}\cdot \mathrm{ev}(v)$ where $\cdot$ is the standard scalar product of $\mathbb{F}^D$. \\

This implies that $H_S = \mathrm{ev}(S)^{\perp} =\ker(\mathrm{ev}(S))$,  so $|H_S| = 2^{D - r_S}$. \\

Let $X(F)$ be the variety of $F \in (Q_d)^n$, denoting by $\varphi(S) := | \{F \in (Q_d)^n \mid S \subseteq X(F) \} |$, we have \[ \varphi(S) = |H_S|^n = 2^{n(D - r_S)} = N^{D - r_S}. \] Indeed the space of system vanishing at least on $S$ is $(H_S)^n$, since a system vanishes on $S$ if and only if all its polynomials vanish on $S$. \\

For $k = 0, \dots, 2^n$, let $\alpha_k := |\{F \in (Q_d)^n \mid |X(F)| = k\}|$ denote the number of square systems with $k$ solutions.
Recalling that $r_V = D$, we show

\begin{theorem}\label{thm:boolean-matroid-link}
In the evaluation matroid $M = (V, r)$, for any $v \in V$:
\[
\alpha_0 = \chi_M(N) \quad \text{and} \quad \alpha_1 = \chi_{M/v}(N)N.
\]
\end{theorem}

\begin{proof}
Applying inclusion-exclusion directly to the complement of the arrangement of all the $(H_v)^n$ inside $\mathbb{F}^{nD}$:
\[
\alpha_0 = \left| \mathbb{F}^{nD} \setminus \bigcup_{v \in V} H_v^n \right| = \sum_{S \subseteq V} (-1)^{|S|} \left| \bigcap_{v \in S} H_v^n \right| = \sum_{S \subseteq V} (-1)^{|S|} |H_S^n| = \sum_{S \subseteq V} (-1)^{|S|} N^{D - r(S)} = \chi_M(N).
\]
For the second identity, fix $v \in V$. Systems whose unique solution is $v$ correspond to $H_v^n \setminus \bigcup_{w \neq v} (H_v^n \cap H_w^n)$. By inclusion-exclusion, their cardinality is:
\begin{align*}
    \left| (H_v)^n \setminus \bigcup_{w \neq v} ((H_v)^n \cap (H_w)^n)  \right| = \sum_{S \subseteq V \setminus \{v\}} (-1)^{|S|} \left| \bigcap_{w \in S} ((H_v)^n \cap (H_w)^n) \right|  \\ =  \sum_{S \subseteq V \setminus \{v\}} (-1)^{|S|} \varphi(S \cup \{v\}) = \sum_{S \subseteq V \setminus \{v\}} (-1)^{|S|}N^{D - r_{S \cup \{v\}}} = \chi_{M/v}(N). 
\end{align*}

Furthermore, the affine transformation group of $V$ acts transitively on its $N$ points, preserves degree, induces an invertible linear transformation of the coefficient space $\mathbb{F}^D$, and thus preserves the count for every point, yielding $\alpha_1 =  \chi_{M/v}(N)N$.
\end{proof}

\section{A sign-reversing involution and the matroid inequality} \label{sec:reduction}

In this section $M=(E,r)$ is an arbitrary matroid and $e\in E$ is neither a loop nor a coloop and has no parallel element. We reduce $\lambda\chi_{M/e}(\lambda)-\chi_M(\lambda)$ to sums over two families of \emph{independent} sets, and then bound it from below.

The three steps are as follows. In \S\ref{subsec:ports} the difference is split
along the \emph{port} of $M$ at $e$: the sets $S$ with $e\in\cl(S)$ contribute one
sum, the remaining ones another. Neither sum is alternating in a useful way, because
the ranks $r(S)$ vary with $S$. In \S\ref{subsec:involutions} we remove that
obstacle with a sign-reversing involution in the spirit of Whitney's
broken-circuit theorem: on each family it pairs off all but the sets carrying no
\emph{active} element, and paired sets have equal rank and opposite sign, so they
cancel. What survives is indexed by \emph{independent} sets only, on which
$r(S)=|S|$, so both sums become polynomials in $\lambda$ with explicit integer
coefficients $a_k$ and $b_k$ counting fixed sets of each size. In
\S\ref{subsec:fixed} we bound those coefficients: each family is closed under
passing to subsets, up to the members of the port itself, which gives
$(k+1)a_{k+1}\le(|E|-1-k)a_k$ and a matching inequality for the $b_k$ with an error
term $\mu_{k+1}$ counting the elements of the port of size $k+1$. Pairing consecutive terms then makes the
first sum non-negative and the second small, which is
Theorem~\ref{thm:matroid-inequality}.

\subsection{Matroid ports}\label{subsec:ports}

Lehman firstly introduced the port of a matroid in 1964 in a work on Shannon's Switching Game \cite{port}.
\begin{definition}
    Let $e \in E$, then the $e$-\emph{port} of $M$, or the \emph{port of M at} $e$, is the set of circuits containing $e$, without $e$. \[P_e := \{ C \setminus \{e\} \mid C \in \mathcal{C}(M), \ e \in C\}\]
    Elements of $P_e$ are called \emph{minimals of} $e$.  
\end{definition}

We define the upfilter generated by it and the associated polynomial.
\begin{definition}
For $e \in E$, define the $e$-\emph{port upfilter} as \[T_e := \{ S \subseteq E \setminus \{e\} \mid \exists  \ C \in P_e \text{ s.t. } C \subseteq S\} \] and the $e$-\emph{port polynomial}
\[
W_e(\lambda) = \sum_{S \in T_e} (-1)^{|S|} \lambda^{r(M) - r(S)}.
\]
\end{definition}

\begin{observation} \label{obs:Te-cl} 
    We observe that $T_e = \{ S \subseteq E \setminus \{e\} \mid e \in \mathrm{cl}(S)\}$. Indeed $e \in \mathrm{cl}(S)$ iff there exists a circuit $C \subseteq S \cup \{e\}$, with $e \in C$.
\end{observation}

\begin{proposition}\label{We}
Let $M$ be a matroid on ground set $E$. Then for every non-loop $e \in E$,
\begin{equation*}
    \chi_{M/e}(\lambda)(\lambda - 1)= \chi_M(\lambda) + (\lambda - 1) W_e(\lambda),
\end{equation*} equivalently
\begin{equation}\label{eq:We}
  \chi_{M/e}(\lambda)\lambda - \chi_M(\lambda) = \chi_{M/e}(\lambda) + (\lambda - 1) W_e(\lambda).    
\end{equation}
\end{proposition}

\begin{proof}
We prove the second equality applying the formula from Proposition \ref{identity}. \\
For any $S \subseteq E \setminus \{e\}$, consider the term differences between $ \chi_{M/e}(\lambda)\lambda$ and $\chi_{M \setminus e}(\lambda)$, i.e. $(-1)^{|S|}(\lambda^{r(M)-r(S \cup \{e\}) + 1} - \lambda^{r(M)-r(S)}) $. If $e \notin \mathrm{cl}(S)$, adjoining $e$ increases rank by 1, yielding matching terms. If $e \in \mathrm{cl}(S)$, the difference evaluates to $(\lambda - 1)(-1)^{|S|} \lambda^{r(M) - r(S)}$.
\begin{align*}
&\chi_{M/e}(\lambda)\lambda - \chi_M(\lambda) \\ &=  \chi_{M/e}(\lambda)\lambda - (\chi_{M\setminus e}(\lambda) - \chi_{M / e}(\lambda)) \\&= \lambda \sum_{S \subseteq E \setminus \{e\}} (-1)^{|S|} \lambda^{r(M) - r(S \cup \{e\})} - \sum_{S \subseteq E \setminus \{e\}} (-1)^{|S|} \lambda^{r(M) - r(S)} + \sum_{S \subseteq E \setminus \{e\}} (-1)^{|S|} \lambda^{r(M) - r(S \cup \{e\})} \\
&= (\lambda - 1) \sum_{\substack{S \subseteq E \setminus \{e\} \\ e \in \mathrm{cl}(S)}} (-1)^{|S|} \lambda^{r(M) - r(S)} + \sum_{S' \subseteq E \setminus \{e\}} (-1)^{|S'|} \lambda^{r(M) - r(S' \cup \{e\})} \\
&= \chi_{M/e}(\lambda) + (\lambda - 1) W_e(\lambda).  \qedhere 
\end{align*}
\end{proof}

This can be translated in the evaluation matroid as

\begin{corollary} \label{cor:We-alpha}
    \begin{equation*}
        \alpha_1 - \alpha_0 = \frac{\alpha_1}{N} + (N - 1)\sum_{S \in T_v}(-1)^{|S|}\varphi(S)
    \end{equation*}
\end{corollary}

\begin{proof}
    Recalling that $\alpha_0 = \chi_{M}(N)$ and $\alpha_1 = \chi_{M/v}(N)N$ (hence $\chi_{M/v} = \frac{\alpha_1}{N}$), and evaluating the formula \eqref{eq:We} in $\lambda = N$, we get the result.
\end{proof}

\subsection{Involutions}\label{subsec:involutions}

The sign of $W_e(\lambda)$ is hard to read off because the ranks of the sets in $T_e$ vary. We use a sign-reversing involution, in the spirit of Whitney's broken-circuit theorem \cite{whitney32} and its bijective proof by Blass and Sagan \cite{blasssagan}, to rewrite each of the sums in \eqref{eq:We} as a sum over independent sets only.

Fix a total order $e_1<e_2<\dots$ on $E$. For $S\subseteq E$ and $p\in E$ let $S_{>p}:=\{s\in S: s>p\}$. Put $E_0 := E$ and $E_i:=E\setminus\{e_1,\dots,e_i\}$, for $i \ge 1$.

\begin{definition}\label{def:active}
Let $i\ge0$ and $S\subseteq E_i$. An element $p\in E_i$ is $i$-\emph{active} for $S$ if $p\in\cl(S_{>p})$. If $S$ has an $i$-active element, its $i$-\emph{pivot} is $p_i(S):=\min\{p\in E_i: p\ i\text{-active for }S\}$ and \\
$\iota_i(S):=S\,\triangle\,\{p_i(S)\}$; otherwise $\iota_i(S):=S$. 
\end{definition}

\begin{definition}
Let $T \subseteq M$, then $\mathrm{Fix}_i(T) := \{S \in T \mid \iota_i(S) = S\}$.
\end{definition}

\begin{observation}
    $p \in \mathrm{cl}(S_{>p})$ if and only if there is a circuit $C \subseteq S_{>p} \cup \{p\}$ with $p \in C$. Moreover $\min C = p$ since $C \setminus \{p\} \subseteq S_{>p}$ and hence exceeds $p$. For any $e_i$ such that $p > e_i$, $C$ is therefore a circuit of $M \setminus \{e_1, \dots, e_i\}$. 
\end{observation}

\begin{lemma}\label{lem:involution-full}
Let $i\ge0$. Then:
\begin{enumerate}[label=\textnormal{(\alph*)}]
\item the map $\iota_i:2^{E_i}\to2^{E_i}$ is an involution, and every non-fixed
orbit consists of two sets $\{S,\iota_i(S)\}$ with $|\iota_i(S)|=|S|\pm1$;
\item if $S\notin\Fix_i(2^{E_i})$ and $p=p_i(S)$, then $p\in\cl(S\setminus\{p\})$;
consequently $r(S\cup X)=r(\iota_i(S)\cup X)$ for every $X\subseteq E$, and in
particular $r(S)=r(\iota_i(S))$;
\item any $S\in\Fix_i(2^{E_i})$ is independent, so $r(S)=|S|$.
\end{enumerate}
\end{lemma}

\begin{proof}
\textbf{Involution}: We prove that for any $S \subseteq M$, $p_i(S) = p_i(\iota_i(S))$. This will give that $\iota_i(\iota_i(S)) = \iota_i(S \triangle \{p_i(S)\}) = (S \triangle \{p_i(S)\}) \triangle \{p_i(S)\} = S$, i.e. $\iota_i$ is an involution. \\
Let $p = p_i(S)$ and $S' = \iota_i(S) = S \triangle \{p\}$. Adding or removing $p$ changes $S_{>q}$ only for $q < p$, and leaves $S_{>q}$ unchanged for $q \ge p$; in particular $S'_{>p} = S_{>p}$, so $p$ remains active for $S'$. \\
We claim no $q < p$ is active for $S'$. For $q < p$ we have $S_{>p} \subseteq S_{>q}$, hence $\mathrm{cl}(S_{>p}) \subseteq \mathrm{cl}(S_{>q})$, and since $p \in \mathrm{cl}(S_{>p})$ we get $p \in \mathrm{cl}(S_{>q})$; therefore
\[
\mathrm{cl}(S_{>q} \cup \{p\}) = \mathrm{cl}(S_{>q}) \quad \text{and} \quad \mathrm{cl}(S_{>q} \setminus \{p\}) \subseteq \mathrm{cl}(S_{>q}).
\]
The equality (when $p$ is added) shows adjoining $p$ to $S_{>q}$ is closure-neutral, so it cannot bring $q$ into the closure; the inclusion (when $p$ is removed) shows deleting $p$ can only shrink the closure. In both directions, $q \in \mathrm{cl}(S'_{>q})$ would force $q \in \mathrm{cl}(S_{>q})$, equivalently $q$ active for $S$, contradicting the minimality of $p$. Hence $p$ is the least active element of $S'$, so $p(S') = p$ and $\iota$ is an involution.

\textbf{Rank preservation and sign reversal}: Write $A = S \setminus \{p\}$, so $\{S, \iota (S)\} = \{A, A \cup \{p\}\}$. The witnessing circuit $C$ with $\min C = p$ and $C \setminus \{p\} \subseteq S_{>p} \subseteq A$ shows $p \in \mathrm{cl}(A)$. Since $|\iota (S)| = |S| \pm 1$, the two members of each non-fixed orbit carry opposite signs and equal rank.

\textbf{Fixed points}: By definition of $\iota_i$, $S$ is fixed iff it has no
$i$-active element. Such an $S$ is independent: if $S$ contained a circuit $C$, then
$p:=\min C$ would satisfy $C\setminus\{p\}\subseteq S_{>p}$, hence
$p\in\cl(S_{>p})$, so $p$ would be $i$-active for $S$, a contradiction. Hence
$r(S)=|S|$, which is (c).
\end{proof}

\begin{corollary}[restricted Whitney theorem]\label{cor:whitney}
If $T\subseteq2^{E_i}$ satisfies $\iota_i(T)=T$ and $X\subseteq E$, then
\[
\sum_{S\in T}(-1)^{|S|}\lambda^{r(M)-r(S\cup X)}=\sum_{S\in\Fix_i(T)}(-1)^{|S|}\lambda^{r(M)-r(S\cup X)} .
\]
In particular, taking $i=0$, $T=2^E$ and $X=\emptyset$,
\[
\chi_M(\lambda)=\sum_{S\in\Fix_0(2^E)}(-1)^{|S|}\lambda^{r(M)-|S|}.
\]
\end{corollary}

\begin{proof}
By Lemma~\ref{lem:involution-full}(a),(b) the non-fixed members of $T$ split into
pairs $\{S,\iota_i(S)\}$ with $|\iota_i(S)|=|S|\pm1$ and
$r(S\cup X)=r(\iota_i(S)\cup X)$, whose contributions cancel. For the second
formula use Lemma~\ref{lem:involution-full}(c) to replace $r(S)$ by $|S|$ on the
fixed sets.
\end{proof}

We now prove the $\iota_i$-invariance of two subsets of $M$, that will appear in the general formula in Corollary \ref{cor:pointed-reduction-full}.

\begin{lemma}\label{lem:invariance-Te-full}
Let $e := e_1$. The following hold:
\begin{itemize}
    \item[1.] $T_e$ is $\iota_1$-invariant and for $S \in \mathrm{Fix}_1(T_e)$, $r(S) = r(S \cup \{e\})=|S|$;
    \item[2.] $\overline{T_e} := \{S \in M \setminus T_e \mid e \notin S\}$ is $\iota_0$-invariant and for $S \in \mathrm{Fix}_0(\overline{T_e})$, $r(S)=|S|$.
\end{itemize}
\end{lemma}

\begin{proof}
\textbf{(1)}: Let $\iota = \iota_1$, $S \in T_e$, $S \ni p = p(S)$ and $A = S \setminus \{p\}$. The closure operator is monotone and idempotent, so $p \in \mathrm{cl}(A)$ forces
\[
\mathrm{cl}(A) \subseteq \mathrm{cl}(A \cup \{p\}) \subseteq \mathrm{cl}(\mathrm{cl}(A)) = \mathrm{cl}(A),
\]
that is, $\mathrm{cl}(A \cup \{p\}) = \mathrm{cl}(A)$. Thus $A$ and $A \cup \{p\}$ have the same closure $\mathrm{cl}(A)$; in particular $\mathrm{cl}(\iota(S)) = \mathrm{cl}(S) = \mathrm{cl}(A)$ and $r(\iota(S)) = r(S)$. Consequently membership of any fixed element in the closure is unaffected by the toggle:
\[
e \in \mathrm{cl}(S) \iff e \in \mathrm{cl}(A) \iff e \in \mathrm{cl}(\iota(S)).
\]
Hence $\iota$ maps $T_e$ to itself. \\
The rank formula is given directly by the definition of $T_e$ and Lemma \ref{lem:involution-full}(c).\\

\textbf{(2)}: Let $S\in\overline{T_e}$. Since $e\notin S$ we have $S_{>e}=S$, so $e\in\cl(S_{>e})$ would mean $S\in T_e$; thus $e$ is not $0$-active for $S$, $p_0(S)\ne e$ if it exists, and $\iota_0(S)\subseteq E\setminus\{e\}$. If $S$ is not fixed, $\cl(\iota_0(S))=\cl(S)\not\ni e$, so $\iota_0(S)\in\overline{T_e}$. Fixed members are independent by Lemma~\ref{lem:involution-full}(c), and $r(S\cup\{e\})=r(S)+1$ since $e\notin\cl(S)$.
\end{proof}

\begin{corollary}\label{cor:pointed-reduction-full}
If $e := e_1$ is not a loop then
\[
 \chi_{M/e}(\lambda)\lambda - \chi_M(\lambda) = \sum_{S \in \mathrm{Fix}_0(\overline{T_e})} (-1)^{|S|} \lambda^{r(M) - r(S \cup \{e\})} + \lambda \sum_{S \in \mathrm{Fix}_1(T_e)} (-1)^{|S|} \lambda^{r(M) - r(S)}.
\]
\end{corollary}

\begin{proof}
Both $\overline{T_e}$ and $T_e$ consist of subsets of $E\setminus\{e\}$, and
$2^{E\setminus\{e\}}=\overline{T_e}\sqcup T_e$. Splitting the contraction formula
of Remark~\ref{rem:facts} accordingly, and using $r(S\cup\{e\})=r(S)$ for
$S\in T_e$ (by definition $e\in\cl(S)$ there), gives
\[
\chi_{M/e}(\lambda)=\underbrace{\sum_{S\in\overline{T_e}}
(-1)^{|S|}\lambda^{r(M)-r(S\cup\{e\})}}_{=:A(\lambda)}
\;+\;\underbrace{\sum_{S\in T_e}
(-1)^{|S|}\lambda^{r(M)-r(S)}}_{=\,W_e(\lambda)} .
\]
By Proposition~\ref{We},
\[
\lambda\chi_{M/e}(\lambda)-\chi_M(\lambda)
=\chi_{M/e}(\lambda)+(\lambda-1)W_e(\lambda)
=\bigl(A(\lambda)+W_e(\lambda)\bigr)+(\lambda-1)W_e(\lambda)
=A(\lambda)+\lambda\,W_e(\lambda).
\]
Finally, $\overline{T_e}$ is $\iota_0$-invariant and $T_e$ is $\iota_1$-invariant
by Lemma~\ref{lem:invariance-Te-full}, so Corollary~\ref{cor:whitney} applies to
each sum: with $(i,T,X)=(0,\overline{T_e},\{e\})$ it replaces $A(\lambda)$ by the
sum over $\Fix_0(\overline{T_e})$, and with $(i,T,X)=(1,T_e,\varnothing)$ it
replaces $W_e(\lambda)$ by the sum over $\Fix_1(T_e)$. This is the stated formula.
\end{proof}

As usual, this translates as
\begin{corollary}
    \[
    \alpha_1 - \alpha_0 = \sum_{S \in \mathrm{Fix}_0(\overline{T_v})} (-1)^{|S|}\varphi(S) + N\sum_{S \in \mathrm{Fix}_1(T_v)} (-1)^{|S|}\varphi(S).
    \]
\end{corollary}

\subsection{The structure of the fixed families and the main inequality}\label{subsec:fixed}

\begin{lemma}\label{lem:interval}
Let $S\in\Fix_1(T_e)$. Then $S\cup\{e\}$ contains exactly one circuit $C$, and $e\in C$; the set $S^e:=C\setminus\{e\}$ is the unique member of $P_e$ contained in $S$; and every $S'$ with $S^e\subseteq S'\subseteq S$ belongs to $\Fix_1(T_e)$.
\end{lemma}

\begin{proof}
$S$ is independent and $e\in\cl(S)\setminus S$, so Remark~\ref{rem:facts}(d) gives the unique circuit $C\ni e$. If $P\in P_e$ and $P\subseteq S$, then $P\cup\{e\}$ is a circuit in $S\cup\{e\}$, hence equals $C$. Let $S^e\subseteq S'\subseteq S$. Then $e\in\cl(S^e)\subseteq\cl(S')$, so $S'\in T_e$; and if $p\in E_1$ were $1$-active for $S'$ then $p\in\cl(S'_{>p})\subseteq\cl(S_{>p})$ would be $1$-active for $S$. So $S'\in\Fix_1(T_e)$.
\end{proof}

For $k\ge0$ set
\[
a_k:=\bigl|\{S\in\Fix_0(\overline{T_e}):|S|=k\}\bigr|,\quad
b_k:=\bigl|\{S\in\Fix_1(T_e):|S|=k\}\bigr|,\quad
\mu_k:=\bigl|\{P\in P_e:|P|=k\}\bigr| .
\]
Then $a_0=1$ ($\emptyset\in\Fix_0(\overline{T_e})$ since $e$ is not a loop) and $b_0=b_1=0$ ($e$ is not a loop and has no parallel element). Corollary~\ref{cor:pointed-reduction-full} reads
\begin{equation}\label{eq:ab}
\lambda\chi_{M/e}(\lambda)-\chi_M(\lambda)=\mathcal S_0(\lambda)+\lambda\,\mathcal S_1(\lambda),
\end{equation}
where
\[
\mathcal S_0(\lambda):=\sum_{k\ge0}(-1)^ka_k\lambda^{r(M)-k-1},\qquad
\mathcal S_1(\lambda):=\sum_{k\ge0}(-1)^kb_k\lambda^{r(M)-k},
\]
and $\mathcal S_1(\lambda)=W_e(\lambda)$, $\mathcal S_0(\lambda)=\chi_{M/e}(\lambda)-W_e(\lambda)$ by the proof of Corollary~\ref{cor:pointed-reduction-full}.

\begin{lemma}[shadow bounds]\label{lem:shadow}
For every $k\ge0$:
\begin{enumerate}[label=\textnormal{(\roman*)}]
\item $(k+1)\,a_{k+1}\le(|E|-1-k)\,a_k$;
\item $b_{k+1}\le(|E|-1-k)\,b_k+\mu_{k+1}$.
\end{enumerate}
\end{lemma}

\begin{proof}
(i) $\Fix_0(\overline{T_e})$ is closed under taking subsets: if $S'\subseteq S$ then $e\notin\cl(S')$ because $\cl(S')\subseteq\cl(S)$, and a $0$-active element of $S'$ is $0$-active for $S$. Count the pairs $(T,S)$ with $S\in\Fix_0(\overline{T_e})$, $|S|=k+1$, $T\subseteq S$, $|T|=k$: each $S$ gives $k+1$ pairs, each $T$ at most $|E|-1-k$ (the supersets of $T$ of size $k+1$ in $E\setminus\{e\}$).

(ii) Let $S\in\Fix_1(T_e)$, $|S|=k+1$. If $S\notin P_e$, then $S^e\subsetneq S$ and for $x\in S\setminus S^e$ the set $S\setminus\{x\}$ lies in $\Fix_1(T_e)$ by Lemma~\ref{lem:interval}. Hence every such $S$ has a subset of size $k$ in $\Fix_1(T_e)$, each of which has at most $|E|-1-k$ supersets of size $k+1$; the members of $\Fix_1(T_e)$ of size $k+1$ lying in $P_e$ number at most $\mu_{k+1}$.
\end{proof}

\begin{theorem}\label{thm:matroid-inequality}
Let $M=(E,r)$ be a matroid and $e\in E$ be neither a loop nor a coloop and have no parallel element. Then, with the notation above,
\begin{enumerate}[label=\textnormal{(\alph*)}]
\item $\mathcal S_0(|E|)=\sum_{k\ \mathrm{even}}|E|^{\,r(M)-k-2}\bigl(|E|a_k-a_{k+1}\bigr) \ge 0$;
\item $|E|\,\mathcal S_1(|E|)\ge-\sum_{j\ \mathrm{odd}}\mu_j\,|E|^{\,r(M)+1-j}$;
\item if every member of $P_e$ has odd cardinality, then $\mathcal S_1(|E|)\le0$.
\end{enumerate}
Consequently
\[
|E|\,\chi_{M/e}(|E|)-\chi_M(|E|)\ \ge\ \mathcal S_0(|E|)-\sum_{j\ \mathrm{odd}}\mu_j\,|E|^{\,r(M)+1-j},
\]
and under the hypothesis of \textnormal{(c)} also $|E|\chi_{M/e}(|E|)-\chi_M(|E|)\le\chi_{M/e}(|E|)$ and $\mathcal S_0(|E|)\ge\chi_{M/e}(|E|)$.
\end{theorem}

\begin{proof}
Write $\lambda=|E|$. \\
(a) Group the terms of $\mathcal S_0$ in pairs $(k,k+1)$ with $k$ even; by Lemma~\ref{lem:shadow}(i), $a_{k+1}\le\lambda a_k$. \\ 
(b) Group the terms of $\mathcal S_1$ in pairs $(k,k+1)$ with $k$ even: $b_k\lambda^{r(M)-k}-b_{k+1}\lambda^{r(M)-k-1}=\lambda^{r(M)-k-1}(\lambda b_k-b_{k+1})\ge-\mu_{k+1}\lambda^{r(M)-k-1}$ by Lemma~\ref{lem:shadow}(ii); multiply by $\lambda$ and sum, writing $j=k+1$. \\ 
(c) Group $\mathcal S_1$ in pairs $(k,k+1)$ with $k$ odd; then $\mu_{k+1}=0$, so $b_{k+1}\le\lambda b_k$ and each pair $-\lambda^{r(M)-k-1}(\lambda b_k-b_{k+1})$ is $\le0$. The consequences follow from \eqref{eq:ab}, from $\lambda\chi_{M/e}-\chi_M=\chi_{M/e}+(\lambda-1)\mathcal S_1$, and from $\mathcal S_0=\chi_{M/e}-\mathcal S_1$.
\end{proof}

\section{The port of the evaluation matroid}\label{subsec:rm-codes}

We now specialise Theorem~\ref{thm:matroid-inequality} to $M=(V,\overline r)$ and
identify the port $P_v$ in coding-theoretic terms. The identification is the point
at which Reed--Muller codes enter: the minimal sets of the port turn out to be the
minimal-support words of the dual code, so the arithmetic quantities $\mu_j$
governing the error term become weight-enumerator coefficients.

\begin{corollary}\label{cor:alpha-inequality}
In the evaluation matroid with $d\le n-1$, for any $v\in V$ and with $a_k,b_k,\mu_k$ computed at $e=v$:
\begin{enumerate}[label=\textnormal{(\alph*)}]
\item $\alpha_1-\alpha_0\ \ge\ \mathcal S_0(N)-\sum_{j\ \mathrm{odd}}\mu_jN^{D+1-j}$, and $\mathcal S_0(N)\ge N^{D-2}(N-a_1)$;
\item if every member of $P_v$ has odd cardinality, then $\alpha_1-\alpha_0\le\alpha_1/N$ and $\alpha_1/N\le\mathcal S_0(N)$.
\end{enumerate}
\end{corollary}

\begin{proof}
Theorem~\ref{thm:matroid-inequality} with $|E|=N$, $r(M)=D$, $\alpha_0=\chi_M(N)$, $\alpha_1/N=\chi_{M/v}(N)$.
\end{proof}

To do this, we connect our matroid construction to coding theory via Reed-Muller codes, which have been introduced by Muller in \cite{RM}. \\
We recall here some basic definition and properties, more can be found at \cite{code}. 

\begin{definition}[\cite{reed,RM}]
For $0 \le d \le n$, the \emph{binary Reed-Muller code} of order $d$ and length $N = 2^n$, denoted $\mathrm{RM}(d,n)$, is the subspace of $\FF^N$ obtained by evaluating all multilinear polynomials in $Q_d$ at all points $v \in V = \mathbb{F}^n$:
\[
\mathrm{RM}(d,n) := \{ (f(v))_{v \in V}  \mid f \in Q_d \}.
\]
In particular, for quadratic systems, $\mathrm{RM}(2,n) = \{ (f(v))_{v \in V} \mid \mathrm{deg}(f) \le 2 \}$.
\end{definition}

\begin{proposition} \label{prop:rm-facts}
    The minimum distance of $RM(d,n)$ is $2^{n - d}$.
\end{proposition}

\begin{definition}
For a linear code $\mathbf{C} \subseteq \mathbb{F}^\nu$, its \emph{dual code} $\mathbf{C}^\perp$ is defined with respect to the standard dot product:
\[
\mathbf{C}^\perp := \{ c \in \mathbb{F}^\nu \mid \forall c' \in \mathbf{C}, \quad \sum_{v \in V} c_v c'_v = 0 \}.
\]
The \emph{support} of a codeword $c \in \mathbb{F}^\nu$ is $\mathrm{supp}(c) := \{ v \in V \mid c_v = 1 \}$, and its Hamming weight is $|c| = |\mathrm{supp}(c)|$.
$c$ is a \emph{minimal support codeword} \cite{ashikhminbarg,minimal} if there is no other word $c' \neq 0$ such that $\mathrm{supp}(c') \subseteq \mathrm{supp}(c)$.
\end{definition}

\begin{proposition}
The dual of the Reed-Muller code $\mathrm{RM}(d,n)$ is another Reed-Muller code given by:
\[
\mathrm{RM}(d,n)^\perp = \mathrm{RM}(n - d - 1, n).
\]
In particular, for quadratic systems, $\mathrm{RM}(2,n)^\perp = \mathrm{RM}(n-3, n)$.
\end{proposition}

\begin{theorem}\label{thm:circuits}
A set $C \subseteq V$ is a circuit of the Boolean evaluation matroid $M = (V, r)$ if and only if $C = \mathrm{supp}(c)$ for a non-zero codeword $c \in \mathrm{RM}(d,n)^\perp$ of minimal support.
\end{theorem}

\begin{proof}
Let $c \in \mathbb{F}^n$. By definition, $c \in \mathrm{RM}(d,n)^\perp$ if and only if $\sum_{v \in V} c_v f(v) = 0$ for all $f \in Q_d$. Since $Q_d$ is spanned by the monomial basis $M_d$, this condition holds if and only if $\sum_{v \in V} c_v m(v) = 0$ for all $m \in M_d$. Under the evaluation map $\mathrm{ev}(v) = (m(v))_{m \in M_d} \in \mathbb{F}^D$, this is equivalent to:
\[
\sum_{v \in V} c_v \mathrm{ev}(v) = \mathbf{0} \quad \text{in } \mathbb{F}^D.
\]
Thus, $c \in \mathrm{RM}(d,n)^\perp$ if and only if the evaluation vectors $\{\mathrm{ev}(v) \mid v \in \mathrm{supp}(c)\}$ are linearly dependent over $\mathbb{F}$.

A circuit $C$ of $M$ is defined as a minimal dependent set of evaluation vectors. Therefore, $C$ is a circuit of $M$ if and only if $C = \mathrm{supp}(c)$ for a non-zero codeword $c \in \mathrm{RM}(d,n)^\perp$ whose support is minimal with respect to set inclusion.
\end{proof}

\begin{corollary}\label{cor:minimal-Tv-RM}
A set $S \subseteq V \setminus \{v\}$ is an element of the $v$-port of $M$ if and only if $S \cup \{v\} = \mathrm{supp}(c)$ for some minimal support codeword $c \in \mathrm{RM}(n - d - 1, n)$.
\end{corollary}

\begin{proof}
This follows directly from Theorem \ref{thm:circuits}.
\end{proof}

\begin{corollary}\label{cor:port-rm}
$P_v=\{\supp(c)\setminus\{v\}: c\in C\text{ minimal-support},\ v\in\supp(c)\}$. Every member of $P_v$ has odd cardinality $\ge2^{d+1}-1$, and
\[
\mu_j=\frac{(j+1)\,A^{\min}_{j+1}(\mathbf{C})}{N}\ \le\ \frac{(j+1)\,A_{j+1}(\mathbf{C})}{N}\qquad(j\ge0).
\]
\end{corollary}

\begin{proof}
The first statement is Theorem~\ref{thm:circuits}; the parity and size statements are Proposition~\ref{prop:rm-facts}. The affine group of $V$ acts transitively on $V$, maps $\mathbf{C}$ onto itself and preserves minimality of supports; hence every point lies in the same number of minimal-support words of weight $j+1$, and counting incidences gives $N\mu_j=(j+1)A^{\min}_{j+1}(\mathbf{C})$.
\end{proof}


\begin{remark}
    It is an open problem to determine the exact distribution of minimal support codewords of $\mathrm{RM(d,n)}$, with $d > 2$, hence we cannot determine precisely the number of minimal sets in $T_v$.
    A better dissertation about this can be found in \cite{minimal}.
\end{remark}

\medskip\noindent
Corollary~\ref{cor:port-rm} converts the error term of
Theorem~\ref{thm:matroid-inequality}(b) into a weighted enumerator of $\mathbf C$.
Throughout the rest of the paper we assume $d\ge2$ and $d\le n-2$, so that $n\ge4$,
$N\ge16$ and $d(\mathbf{C})=2^{d+1}\ge8$; we fix $v\in V$ and compute
$a_k,b_k,\mu_k$ at $e=v$, and we put
\[
\Phi:=\sum_{w}w\,A_w(\mathbf{C})\,N^{-w}.
\]

\begin{lemma}\label{lem:small}
$a_k=\binom{N-1}k$ for $0\le k\le6$. Consequently
\[
\mathcal S_0(N)\ \ge\ N^{D-2}+N^{D-4}\,\frac{(N-1)(N-2)(2N+3)}{6},
\qquad
\sum_{j\ \mathrm{odd}}\mu_jN^{D+1-j}\ \le\ N^{D+1}\,\Phi ,
\]
and $(N-1)(N-2)(2N+3)/6\ge N^3/4$ for $N\ge16$.
\end{lemma}

\begin{proof}
Circuits have size $\ge d(\mathbf{C})\ge8$. Let $S\subseteq V\setminus\{v\}$, $|S|\le6$. Then $S\cup\{v\}$ contains no circuit, so $v\notin\cl(S)$; and $S$ has no $0$-active element, since a $0$-active $p$ would give a circuit $C'$ with $C'\setminus\{p\}\subseteq S_{>p}\subseteq S$ and $|C'\setminus\{p\}|\ge7$. So $S\in\Fix_0(\overline{T_v})$. The bound on $\mathcal S_0$ is Theorem~\ref{thm:matroid-inequality}(a) with $a_0=1$, $a_1=N-1$, $a_2=\binom{N-1}2$, $a_3=\binom{N-1}3$, using $N\binom{N-1}2-\binom{N-1}3=(N-1)(N-2)(2N+3)/6$. The bound on $\sum_j\mu_jN^{D+1-j}$ is Corollary~\ref{cor:port-rm}: $\sum_{j}(j+1)A_{j+1}N^{D-j}=\sum_ww A_wN^{D+1-w}$. The last claim is $N^3/12-N^2/2-5N/6+1\ge0$ for $N\ge16$.
\end{proof}

\section{Bounding the weighted enumerator}\label{sec:rm-bound}

It remains to bound $\Phi=\sum_w wA_w(\mathbf C)N^{-w}$, a weighted enumerator of
$\mathbf C=\RM(n-d-1,n)$ evaluated near $x=1/N$. Three facts drive the estimate.

First, MacWilliams duality turns $\Phi$ into an average over the \emph{dual} code
$\mathbf C^\perp=\RM(d,n)$ of an explicit function
$f(Y_q)=\kappa e^{\lambda Y_q}(NY_q-1)$ of the normalised weight
$Y_q=1-2w(q)/N\in[-1,1]$ (Lemma~\ref{lem:macwilliams_derivative}).

Second, because $\mathbf C$ has minimum distance $M=2^{d+1}$, the first $M-1$
moments of $Y$ over $\mathbf C^\perp$ coincide with those over the whole of
$\FF^N$ (Lemma~\ref{lem:delsarte_orthogonality}, a Delsarte-type orthogonality),
and the full-space average of any polynomial of degree $<M$ in $Y$ vanishes
(Lemma~\ref{lem:binomial_identity}). Hence the Taylor polynomial of $f$ of order
$M-1$ contributes nothing, and $\Phi$ is exactly the average of the Taylor
remainder.

Third, that remainder is controlled by a single bound on $f^{(M)}$
(Lemma~\ref{lem:global_derivative}) together with a count of the minimum-weight
words of $\mathbf C$ (Lemma~\ref{lem:min_weight}). The outcome is
$\Phi=O(N^{-2^d})$, which for $d=2$ is all that Section~\ref{sec:mainthm} consumes.

The constants below are deliberately crude: only the exponent of $N$ matters
downstream, and no attempt is made to optimise them.

\begin{lemma}[a vanishing binomial sum]
\label{lem:binomial_identity}
For any positive integer $N$,
\[
\sum_{w=0}^N \binom{N}{w} \left(\frac{N-1}{N+1}\right)^w (N - 1 - 2w) = 0
\]
\end{lemma}

\begin{proof}
Let $a = \frac{N-1}{N+1}$. Using $(1+a)^N = \sum \binom{N}{w} a^w$ and $N a (1+a)^{N-1} = \sum w \binom{N}{w} a^w$:
\[
\sum_{w=0}^N \binom{N}{w} a^w (N - 1 - 2w) = (1+a)^{N-1} \left[ (N-1)(1+a) - 2Na \right]
\]
Substituting $1+a = \frac{2N}{N+1}$ gives $(N-1)\left(\frac{2N}{N+1}\right) - 2N\left(\frac{N-1}{N+1}\right) = 0$.
\end{proof}

\begin{lemma}[minimum-weight words of $\mathbf{C}$]
\label{lem:min_weight}
The minimum distance of $\mathbf{C} = RM(n-d-1, n)$ is $M = 2^{d+1}$. The number of minimum-weight codewords satisfies:
\[
A_M \le \frac{N^{d+2}}{2^{d+1} K_d} \quad \text{where } K_d = 2^{d(d+1)/2} \prod_{i=1}^{d+1} (2^i - 1)
\]
\end{lemma}

\begin{proof}
Weight-$M$ codewords are indicator functions of $(d+1)$-dimensional affine subspaces of $\mathbb{F}_2^n$. The number of $(d+1)$-dimensional linear subspaces is given by the Gaussian binomial coefficient:
\[
\genfrac{[}{]}{0pt}{}{n}{d+1}_2 = \frac{\prod_{i=0}^d (N - 2^i)}{\prod_{i=0}^d (2^{d+1} - 2^i)}
\]
Bounding the numerator by $N^{d+1}$ and factoring $2^i$ out of each term in the denominator yields:
\[
\prod_{i=0}^d (2^{d+1} - 2^i) = 2^{\sum_{i=0}^d i} \prod_{i=0}^d (2^{d+1-i} - 1) = 2^{d(d+1)/2} \prod_{j=1}^{d+1} (2^j - 1) = K_d
\]
Thus $\genfrac{[}{]}{0pt}{}{n}{d+1}_2 < \frac{N^{d+1}}{K_d}$. Multiplying by the $N/2^{d+1}$ parallel cosets yields $A_M = \frac{N}{2^{d+1}} \genfrac{[}{]}{0pt}{}{n}{d+1}_2 < \frac{N^{d+2}}{2^{d+1} K_d}$.
\end{proof}

\begin{lemma}[MacWilliams form of $\Phi$]
\label{lem:macwilliams_derivative}
Let $\mathbf{C}^\perp = RM(d, n)$ and  \[ g(w) = \frac{1}{N+1} \left(\frac{N+1}{N}\right)^{N-1} \left(\frac{N-1}{N+1}\right)^{w-1} (N - 1 - 2w). \]Then:
\[
\Phi = \frac{1}{|\mathbf{C}^\perp|} \sum_{q \in \mathbf{C}^\perp} g(w(q))
\]
\end{lemma}

\begin{proof}
Let $W_C(x, y) = \sum_{w=0}^N A_w x^w y^{N-w}$. Differentiating gives $\Phi = \frac{1}{N} \left. \frac{\partial W_C}{\partial x} \right|_{(1/N, 1)}$. Applying MacWilliams duality \cite[Ch.~5]{code} $W_C(x, y) = \frac{1}{|\mathbf{C}^\perp|} \sum_{q \in \mathbf{C}^\perp} (y+x)^{N-w(q)} (y-x)^{w(q)}$ and differentiating term-by-term yields:
\begin{multline*}
    \frac{1}{N} \left. \frac{\partial}{\partial x} \left[ (y+x)^{N-w} (y-x)^w \right] \right|_{(1/N, 1)} = \\ = \frac{1}{N} \left(\frac{N+1}{N}\right)^{N-w-1} \left(\frac{N-1}{N+1}\right)^{w-1}\left(\frac{N}{N+1}\right)^{w-1} (N - 1 - 2w) = \\ = \frac{1}{N+1} \left(\frac{N+1}{N}\right)^{N-1} \left(\frac{N-1}{N+1}\right)^{w-1} (N - 1 - 2w) = g(w)
\end{multline*}


\end{proof}

\begin{lemma}[Delsarte orthogonality \cite{delsarte}]
\label{lem:delsarte_orthogonality}
Let $Y_q = 1 - \frac{2w(q)}{N} \in [-1, 1]$. For every integer $0 \le k \le M-1$ where $M = 2^{d+1}$:
\[
\frac{1}{|\mathbf{C}^\perp|} \sum_{q \in \mathbf{C}^\perp} Y_q^k = \frac{1}{2^N} \sum_{u \in \mathbb{F}_2^N} Y_u^k
\]
\end{lemma}

\begin{proof} Notice that $Y_q = \frac{1}{N}\sum_{x=1}^{N}(-1)^{q_x}$.
Expanding $Y_q^k = \frac{1}{N^k} \sum_{x_1, \dots, x_k = 1}^N (-1)^{q \cdot v}$ for $v = 1_{x_1} + \dots + 1_{x_k}$, where $1_{x_i}$ is the vector of weight 1, with non-zero entry in the $i$-th position. Summing over $\mathbf{C}^\perp$ yields the indicator function of $\mathbf{C}$, evaluated on $v$. Namely it is non-zero if and only if $v \in (\mathbf{C}^\perp)^\perp = \mathbf{C}$. Because $v$ has weight at most $k \le M - 1 < d(\mathbf{C})$, $v \in \mathbf{C} \iff v = \mathbf{0}$. The same holds in the second summation, for $(\mathbb{F}_2^N)^\perp = \{0\}$.
\end{proof}

\begin{lemma}[derivative bound]
\label{lem:global_derivative}
Let \[f(Y_q) = g(w(q)) = \kappa e^{\lambda Y_q} (NY_q - 1),\] where $\lambda = \frac{N}{2} \ln\left(\frac{N+1}{N-1}\right)$ and $\kappa = \frac{1}{N+1} \left(\frac{N+1}{N}\right)^{N-1} \left(\frac{N-1}{N+1}\right)^{N/2-1}$. For any $M = 2^{d+1}$ and all $\xi \in [-1, 1]$:
\[
f^{(M)}(\xi) \le e (M+4)
\]
\end{lemma}

\begin{proof}
Differentiating $f(Y)$ $M$ times via the product rule yields $f^{(M)}(\xi) = \kappa e^{\lambda \xi} L(\xi)$, where $L(\xi) = (N\xi - 1)\lambda^M + M N \lambda^{M-1}$. Moreover:
\[
\frac{d}{d\xi} f^{(M)}(\xi) = f^{(M)+1}(\xi) =\kappa \lambda^M e^{\lambda \xi} \left[ N\lambda \xi + (M+1)N - \lambda \right]
\]
Evaluating the linear factor at its minimum $\xi = -1$ yields $(M+1)N - (N+1)\lambda > 144 - 18.13 > 0$ for all $N \ge 16, M \ge 8$. Thus $\frac{d}{d\xi} f^{(M)}(\xi) > 0$ on $[-1, 1]$, proving $f^{(M)}(\xi)$ is strictly increasing on $[-1, 1]$ and maximized uniquely at $\xi = 1$:
\[
f^{(M)}(\xi) \le f^{(M)}(1) = \kappa e^\lambda \left[ (N - 1)\lambda^M + M N \lambda^{M-1} \right]
\]
For the constant,
\[
\kappa e^\lambda = \frac{1}{N+1}\left(\frac{N+1}{N}\right)^{N-1}\left(\frac{N+1}{N-1}\right) = \frac{1}{N-1}\left(1 + \frac{1}{N}\right)^{N-1} < \frac{e}{N-1}.
\]
For the bracket, factor out $\lambda^{M-1}$ and use $\lambda>1$:
\[
L(1) = (N-1)\lambda^{M} + MN\lambda^{M-1}
 = \lambda^{M-1}\bigl[(N-1)\lambda + MN\bigr] .
\]
Since $\ln\frac{N+1}{N-1}=2\sum_{k\ge0}\frac{1}{(2k+1)N^{2k+1}}$, we have
$\lambda = 1+\frac{1}{3N^{2}}+\frac{1}{5N^{4}}+\dots \le 1+\frac{1}{2N^{2}}$, whence
$\lambda^{M-1}\le\bigl(1+\frac{1}{2N^{2}}\bigr)^{M-1}<e^{(M-1)/(2N^{2})}<e$,
because $M-1<2N^{2}$. Moreover $(N-1)\lambda+MN\le(N-1)(M+4)$: this is
$(N-1)(\lambda-1)\le 3N-M-3$, and the left side is at most $1$ while
$3N-M-3\ge N-3>1$, using $M\le N/2$. Therefore
$L(1) < e\,(N-1)(M+4)$ and
\[
f^{(M)}(1) = \kappa e^{\lambda}L(1) < \frac{e}{N-1}\cdot e\,(N-1)(M+4) = e^{2}(M+4).
\]
\end{proof}

\medskip\noindent The three ingredients now combine.

\begin{theorem}[the bound on $\Phi$]
\label{prop:Phi}
Let $\mathbf{C} = RM(n-d-1, n)$ be a binary Reed-Muller code of length $N = 2^n$ ($n \ge 4$, $2 \le d \le n-2$) with dual $\mathbf{C}^\perp = RM(d, n)$. Let $M = 2^{d+1}$ and $K_d = 2^{d(d+1)/2} \prod_{i=1}^{d+1} (2^i - 1)$. Then the weighted sum $\Phi = \sum_{w=1}^N w A_w N^{-w}$ satisfies the following inequality:
\[
\Phi \le \frac{e^2(M+4)}{M!} \left( \frac{(M-1)!!}{N^{2^d}} + \frac{M!}{2^{d+1} K_d \cdot N^{2^{d+1} - d - 2}} \right)
\]
\end{theorem}

\begin{proof}
Set $M = 2^{d+1}$. Expand $f(Y_q)$ in its $(M-1)$-th degree Taylor polynomial around $Y = 0$:
\[
f(Y_q) = P_{M-1}(Y_q) + \frac{f^{(M)}(\xi_q)}{M!} Y_q^M \quad \text{where } \xi_q \in [-1, 1]
\]
Averaging over $\mathbf{C}^\perp$, Lemma~\ref{lem:delsarte_orthogonality} implies $\frac{1}{|\mathbf{C}^\perp|} \sum_{q \in \mathbf{C}^\perp} P_{M-1}(Y_q) = \frac{1}{2^N}\sum_{u \in \FF^N}P_{M-1}(Y_u)$. Since $\frac{1}{2^N}\sum_{u \in \FF^N}P_{M-1}(Y_u) = 0$ by Lemma~\ref{lem:binomial_identity} and the Taylor remainder $R_M(Y_u) = \frac{f^{(M)}(\xi_u)}{M!} Y_u^M \ge 0$, we have $\Phi = \frac{1}{|\mathbf{C}^\perp|}\sum_{q \in \mathbf{C}^\perp}R_M(Y_q)$. Applying Lemma~\ref{lem:global_derivative} gives:
\[
\Phi \le \frac{e^2(M+4)}{M!} \left( \frac{1}{|\mathbf{C}^\perp|} \sum_{q \in \mathbf{C}^\perp} Y_q^M \right)
\]
Expanding the $M$-th moment into character indicators $v = 1_{x_1} + \dots + 1_{x_M}$:
\[
\frac{1}{|\mathbf{C}^\perp|} \sum_{q \in \mathbf{C}^\perp} Y_q^M = \frac{1}{N^M} \left| \{ (x_1, \dots, x_M) \mid v \in \mathbf{C} \} \right|
\]
Since $\text{wt}(v) \le M = d(\mathbf{C})$, $v \in \mathbf{C}$ if and only if $v = \mathbf{0}$ (contributing $(M-1)!! N^{2^d}$ paired tuples) or $v \in \mathbf{C} \setminus \{\mathbf{0}\}$ (contributing $M! A_M$ weight-$M$ tuples). Substituting Lemma~\ref{lem:min_weight} for $A_M$:
\[
\frac{1}{|\mathbf{C}^\perp|} \sum_{q \in \mathbf{C}^\perp} Y_q^M \le \frac{(M-1)!! N^{2^d} + M! \left( \frac{N^{d+2}}{2^{d+1} K_d} \right)}{N^{2^{d+1}}} = \frac{(M-1)!!}{N^{2^d}} + \frac{M!}{2^{d+1} K_d \cdot N^{2^{d+1} - d - 2}}
\]
Multiplying by $\frac{e^2(M+4)}{M!}$ completes the proof.
\end{proof}

\begin{corollary}
\label{cor:d2_eval}
For any $2 \le d \le n - 2$, $\Phi < \frac{1}{3 N^4}$.
\end{corollary}

\begin{proof}
 Let $b_d$ represent the bound of the previous Theorem, then for any $d \ge 2$, $b_d \le b_2$, since the power of $N$ grows exponentially. \\
 Explicitly computing $b_2$ we have that for $d = 2$, $M = 2^3 = 8$, $K_2 = 2^3 \cdot 21 = 168$, $2^d = 4$, and $2^{d+1} - d - 2 = 4$. Hence it evaluates to:
\[
\Phi \le \frac{12e^2}{40320} \left( \frac{105}{N^4} + \frac{40320}{8 \cdot 168 \cdot N^4} \right) = \frac{12e^2}{40320} \left( \frac{105 + 30}{N^4} \right) = \frac{1620 e^2}{40320 N^4} < \frac{1}{3 N^4}
\]
\end{proof}

\section{Proof of the main theorem}\label{sec:mainthm}

We can now prove Theorem~\ref{thm:intro}, restated here as
Theorem~\ref{thm:main}. The three cases $d=n$, $d=n-1$ and $d\le n-2$ are handled
separately: the first two are direct computations, and only the last needs the
machinery of Sections~\ref{sec:reduction}--\ref{sec:rm-bound}. A final remark shows
that the hypothesis $d\ge2$ cannot be dropped.

\begin{theorem}\label{thm:main}
Let $n\ge2$ and $2\le d\le n$. Then
\begin{enumerate}[label=\textnormal{(\roman*)}]
\item $\alpha_1-\alpha_0\le\alpha_1/N$, with equality if and only if $d=n$;
\item $\alpha_1>\alpha_0$;
\item if $d\le n-2$, then
\[
\frac{\alpha_1}{N}-\frac{N^{D-3}}{3}\;\le\;\alpha_1-\alpha_0\;\le\;\frac{\alpha_1}{N},
\qquad
\frac{\alpha_1}{N}\;\ge\;N^{D-4}\,\frac{(N-1)(N-2)(2N+3)}{6}\;\ge\;\frac{N^{D-1}}4,
\]
and hence $\bigl(1-\frac4{3N^2}\bigr)\frac{\alpha_1}N\le\alpha_1-\alpha_0\le\frac{\alpha_1}N$.
\end{enumerate}
\end{theorem}

\begin{proof}
\emph{Case $d=n$.} $Q_d=Q$ is the space of all functions $V\to\FF$, so a system is an arbitrary map $V\to\FF^n$ and $\alpha_0=(N-1)^N$, $\alpha_1=N(N-1)^{N-1}$ (choose the unique zero and the non-zero values elsewhere; equivalently, the number of solutions is binomial with parameters $N,1/N$, cf.\ \cite[Theorem~1]{jain}). Thus $\alpha_1-\alpha_0=\alpha_1/N>0$.

\emph{Case $d=n-1$ ($n\ge3$).} Here $\mathbf{C}=\RM(0,n)=\{\mathbf0,\mathbf1\}$, so the only circuit is $V$ and $T_v=\{V\setminus\{v\}\}$; moreover $r(V\setminus\{v\})=D$ by Lemma~\ref{lem:coloop}, so $\varphi(V\setminus\{v\})=1$ and $W_v(N)=(-1)^{N-1}=-1$. Corollary~\ref{cor:We-alpha} gives $\alpha_1-\alpha_0=\alpha_1/N-(N-1)<\alpha_1/N$. Also $\alpha_1/N\ge|GL(n,\FF)|>N-1$ (every invertible linear system has $\mathbf0$ as its unique zero, and $|GL(n,\FF)|=\prod_{i<n}(2^n-2^i)>2^n-1$ for $n\ge2$), so $\alpha_1>\alpha_0$.

\emph{Case $d\le n-2$.} Then $n\ge4$, $N\ge16$ and $d(\mathbf{C})=2^{d+1}\ge8$. By Corollary~\ref{cor:port-rm} every member of $P_v$ has odd cardinality, so Corollary~\ref{cor:alpha-inequality}(b) gives
\[
\alpha_1-\alpha_0\le\frac{\alpha_1}N\qquad\text{and}\qquad \mathcal S_0(N)\ge\frac{\alpha_1}N .
\]
By Corollary~\ref{cor:alpha-inequality}(a), Lemma~\ref{lem:small} and Proposition~\ref{prop:Phi},
\[
\alpha_1-\alpha_0\ \ge\ \mathcal S_0(N)-\sum_{j\ \mathrm{odd}}\mu_jN^{D+1-j}\ \ge\ \mathcal S_0(N)-N^{D+1}\Phi\ \ge\ \frac{\alpha_1}N-\frac{N^{D-3}}{3}.
\]
For the lower bound on $\alpha_1/N$, note that $\alpha_1/N=\chi_{M/v}(N)=\mathcal S_0(N)+\mathcal S_1(N)$ by \eqref{eq:ab} and Theorem~\ref{thm:boolean-matroid-link}, and $\mathcal S_1(N)\ge-\frac1N\sum_{j\ \mathrm{odd}}\mu_jN^{D+1-j}\ge-N^{D}\Phi\ge-N^{D-4}/3$ by Theorem~\ref{thm:matroid-inequality}(b), Lemma~\ref{lem:small} and Proposition~\ref{prop:Phi}. Hence, by Lemma~\ref{lem:small},
\begin{align*}
\frac{\alpha_1}N\ &\ge\ N^{D-2}\Bigl(1-\frac1{3N^2}\Bigr)+N^{D-4}\frac{(N-1)(N-2)(2N+3)}6\\
&\ge\ N^{D-4}\frac{(N-1)(N-2)(2N+3)}6\ \ge\ \frac{N^{D-1}}4 .
\end{align*}
Therefore $N^{D-3}/3\le\frac4{3N^2}\cdot\frac{\alpha_1}N$, which proves (iii), and (ii) follows since $\alpha_1>0$.

\emph{Strictness in (i) for $d\le n-2$.} Since $\alpha_1-\alpha_0=\alpha_1/N+(N-1)W_v(N)$ (Corollary~\ref{cor:We-alpha}), it suffices to show $W_v(N)<0$. The value $W_v(N)=\mathcal S_1(N)$ does not depend on the order chosen on $V$, so we may choose it: let $P_0\in P_v$ have the minimum cardinality $m$ (odd) and order $V$ so that $v$ comes first and the elements of $P_0$ next. Then $P_0\in\Fix_1(T_v)$: an element $p\in P_0$ cannot be $1$-active since $P_0$ is independent, and an element $p\notin P_0\cup\{v\}$ has $(P_0)_{>p}=\emptyset$. Hence $b_m\ge1$, and in the proof of Theorem~\ref{thm:matroid-inequality}(c) the pair $(m,m+1)$ contributes $-N^{D-m-1}(Nb_m-b_{m+1})\le-N^{D-m-1}(N-(N-1-m))b_m<0$, all other pairs being $\le0$. So $\mathcal S_1(N)<0$.
\end{proof}

\begin{remark}\label{rem:final}
The hypothesis $d\ge2$ cannot be dropped. For $d=1$ a system is an affine map
$S(x)=Ax+b$ with $A\in\FF^{n\times n}$ and $b\in\FF^n$, so $D=n+1$ and the $N^{D}$
systems correspond to the pairs $(A,b)$. The set $S^{-1}(0)$ is empty when
$b\notin\Im A$ and is otherwise a coset of $\ker A$, of size $2^{\,n-\mathrm{rank}\,A}$; hence,
writing $g_r$ for the number of $A$ of rank $r$,
\[
\alpha_0=\sum_{r=0}^{n}g_r\bigl(N-2^{r}\bigr),
\qquad
\alpha_1=N\,g_n=N\,|\mathrm{GL}(n,\FF)| ,
\]
since a unique solution means $\mathrm{rank}\,A=n$, and then every $b$ occurs. Now
$g_r=\genfrac{[}{]}{0pt}{}{n}{r}_2\prod_{i=0}^{r-1}(N-2^{i})$, and writing
$P=\prod_{i=0}^{n-3}(N-2^{i})$ one has
\[
g_n=P\Bigl(\tfrac{3N}4\Bigr)\Bigl(\tfrac N2\Bigr),\qquad
g_{n-1}=(N-1)P\Bigl(\tfrac{3N}4\Bigr),\qquad
g_{n-2}=\tfrac{(N-1)(\frac N2-1)}3\,P .
\]
Keeping only the terms $r=n-1$ and $r=n-2$ of $\alpha_0$ (all terms are
non-negative, and the term $r=n$ vanishes) gives
\[
\alpha_0\ \ge\ g_{n-1}\tfrac N2+g_{n-2}\tfrac{3N}4
=(N-1)P\Bigl(\tfrac{N^{2}}2-\tfrac N4\Bigr),
\qquad
\alpha_1=\tfrac38PN^{3},
\]
so that
\[
\frac{\alpha_0}{\alpha_1}\ \ge\ \frac{4N^{2}-6N+2}{3N^{2}}\ >\ 1
\qquad\Longleftrightarrow\qquad N^{2}-6N+2>0 ,
\]
which holds for $N\ge8$. Hence $\alpha_0>\alpha_1$ for every $n\ge3$, the reverse
of the inequality of Theorem~\ref{thm:main}. \textup{(}For $n=1,2$ one computes
$(\alpha_0,\alpha_1)=(1,2)$ and $(21,24)$, so $\alpha_0<\alpha_1$ there; the ratio
$\alpha_0/\alpha_1$ increases to $4/3$.\textup{)}
\end{remark}

\begin{proof}[Proof of \eqref{eq:intro-mq}]
Apply Theorem~\ref{thm:main} with $d=2$ and $n\ge2$, and use $\alpha_1-\alpha_0\le\alpha_1/N\iff\alpha_1\le(1+\frac1{N-1})\alpha_0$. Equality on the right holds iff $d=n$, i.e.\ iff $n=2$.
\end{proof}


\section{Remarks on complexity}\label{sec:complexity}

We record the elementary facts used in the introduction. A \emph{quadratic system} over $\FF$ is a tuple $F=(f_1,\dots,f_m)$ with $f_i\in Q_2$ (in some number $n$ of variables), and $X(F)\subseteq\FF^n$ is its set of solutions. A polynomial-time map $F\mapsto F'$ between systems is \emph{parsimonious} if $|X(F')|=|X(F)|$.

\begin{proposition}[squaring]\label{prop:square}
Let $F$ be a quadratic system over $\FF$ in $n$ variables with $m$ equations.
\begin{enumerate}[label=\textnormal{(\alph*)}]
\item If $m\le n$, appending $n-m$ equations $0=0$ gives an $n\times n$ system with
the same solution set.
\item If $m>n$, adjoining $m-n$ variables not occurring in any equation gives an
$m\times m$ system whose solution set is $X(F)\times\FF^{\,m-n}$; this is square but
not parsimonious, the number of solutions being multiplied by $2^{\,m-n}$.
\item If $m>n$, replacing a pair of equations $f=0$, $g=0$ by
\[
w+f=0,\qquad u+g=0,\qquad w+u+wu=0
\]
in two new variables $w,u$, and iterating $m-n$ times, gives a
$(2m-n)\times(2m-n)$ system with the same \emph{number} of solutions.
\end{enumerate}
All three maps are computable in polynomial time.
\end{proposition}

\begin{proof}
Only (c) needs a word. For each $x$ the first two equations force $w=f(x)$ and
$u=g(x)$, and $w+u+wu=0$ holds iff $w=u=0$; so the solutions of the new system are
in bijection with those of the old one, and the equations are in $Q_2$. One step
replaces two equations by three and adds two variables, so it leaves the number of
equations minus the number of variables one lower than before: starting from the
surplus $m-n>0$, after $m-n$ steps the system has $m+(m-n)=2m-n$ equations in
$n+2(m-n)=2m-n$ variables, hence is square.
\end{proof}

\begin{proposition}[clauses]\label{prop:clauses}
There is a parsimonious polynomial-time map from CNF formulas to quadratic systems over $\FF$.
\end{proposition}

\begin{proof}
Identify a literal $\ell$ in the variables $x_1,\dots,x_n$ with the affine function $x_i$ or $1+x_i$, so that $\ell$ is true iff the function equals $1$. For a clause $C=\ell_1\vee\dots\vee\ell_k$, $k\ge2$, introduce variables $z_1,\dots,z_{k-2}$ and the equations
\[
z_1+\ell_1+\ell_2+\ell_1\ell_2=0,\qquad z_j+z_{j-1}+\ell_{j+1}+z_{j-1}\ell_{j+1}=0\ \ (2\le j\le k-2),\qquad z_{k-2}+\ell_k+z_{k-2}\ell_k+1=0
\]
(for $k=2$ the single equation $\ell_1+\ell_2+\ell_1\ell_2+1=0$; for $k=1$ the equation $\ell_1+1=0$). Since $a\vee b$ is the function $a+b+ab$, the first $k-2$ equations force $z_j=\ell_1\vee\dots\vee\ell_{j+1}$ for every assignment of the $x_i$, and the last one holds iff $C$ is satisfied. The equations are in $Q_2$, and an assignment of the $x_i$ extends to a solution of the system in exactly one way iff it satisfies $C$, in no way otherwise. Doing this for every clause, with disjoint sets of auxiliary variables, gives the map.
\end{proof}

\begin{corollary}\label{cor:complexity}
Let $\MQ_0(n)$, $\MQ_1(n)$ be as in the introduction.
\begin{enumerate}[label=\textnormal{(\alph*)}]
\item Deciding whether a square quadratic system over $\FF$ has a solution is NP-complete; hence $\bigcup_{n\ge2}\MQ_0(n)$ is coNP-complete.
\item $\bigcup_{n\ge2}\MQ_1(n)$ is in DP, is coNP-hard, and is NP-hard under randomized polynomial-time reductions.
\end{enumerate}
\end{corollary}

\begin{proof}
Propositions~\ref{prop:clauses} and~\ref{prop:square} compose to a parsimonious polynomial-time reduction $\varphi\mapsto F_\varphi$ from CNF-SAT to square quadratic systems: $\varphi$ has exactly $s$ satisfying assignments iff $F_\varphi$ has exactly $s$ solutions. (a) Membership in NP is clear, hardness follows from the reduction (or from \cite{fraenkelyesha}), and the complement of an NP-complete language is coNP-complete. (b) A square system has exactly one solution iff it has at least one (an NP property) and at most one (a coNP property: two distinct solutions are a certificate of failure); so the language is in DP. The reduction maps the unique-satisfiability language $\{\varphi:\varphi\text{ has exactly one satisfying assignment}\}$ into $\bigcup_n\MQ_1(n)$ and its complement into the complement, so it is a many-one reduction from unique satisfiability, which is coNP-hard \cite{blassgurevich} and NP-hard under randomized reductions \cite{valiantvazirani}.
\end{proof}

By Theorem~\ref{thm:main}, for every $n\ge2$ there exist injections $\MQ_0(n)\hookrightarrow\MQ_1(n)$ whose images omit at most a $2^{-n}$ fraction of $\MQ_1(n)$. Nothing in this paper says how to compute such an injection; an injection computable in polynomial time, together with a polynomial-time inverse on its image, would give a rather unusual kind of reduction between a coNP-complete language and a language in DP. We leave this as an open question.

\section*{Acknowledgements}

These results are included in the first author's MSC thesis, supervised by the
second author. Several of the numerical claims were checked by computer algebra
with the assistance of Anthropic Claude and OpenAI ChatGPT, which also reviewed the
arguments and assisted with some proofs.

\newpage
\bibliographystyle{plain}
\bibliography{biblio}

\end{document}